\section{Morphisms as operators}\label{chap:operator}

The morphisms used in Section~\ref{chap:functoriality} were the maps of base
spaces $\underline{F}:Y\to X$ themselves.
Under that definition a $\operatorname{Hom}$ set is no more than a subset of
$\operatorname{Map}(Y,X)$, so the uniqueness of the factorization in
Theorem~\ref{thm:main-theorem} followed automatically from the agreement of the
maps of base spaces.
That is, for the non-ae version of the main theorem the uniqueness part has no
content (this is not so for the ae version; see
Subsection~\ref{sec:main-theorem-ae}).

In Section~\ref{chap:operator} we detach morphisms from maps of base spaces and
redefine them as linear maps between Riesz spaces.

\begin{quote}
  \textbf{A note on attribution.}
  This formulation is based on the one introduced in \cite{Yamagami} in the
  context of discussing iterated integrals, but it is \textbf{not the same
  thing.}
  That book imposes \emph{order continuity} in addition to linearity and
  positivity, whereas the operator morphisms of integration systems in
  Section~\ref{chap:operator} (Definition~\ref{def:operator-morphism}) do not
  impose it.
  That is, what we adopt is that book's definition with one condition removed.

  Order continuity is a condition that book imposes in order to
  \emph{construct} $J\circ\Lambda$ as an integral, and it becomes unnecessary in
  the setting of Section~\ref{chap:operator}, which requires preservation of the
  integral.
  For details see Remark~\ref{rem:order-continuity-redundant}.
\end{quote}

\noindent
Note that the definition of \cite{Yamagami}, with order continuity imposed
(Remark~\ref{rem:yamagami-operator}), includes Markov operators and transition
probabilities, and that book says so explicitly.

This replacement changes the following three things.

\begin{itemize}
  \item Morphisms no longer pass through the base space, and a
    $\operatorname{Hom}$ set is no longer a subset of a set of maps.
  \item Uniqueness of the factorization no longer follows automatically; it is
    obtained using the density of $L$ in $L^1$
    (Lemma~\ref{lem:l1-approx-by-source}).
    The same density is used in the ae version in
    Subsection~\ref{sec:main-theorem-ae}.
  \item The extension is determined only up to null functions.
    It therefore becomes natural to divide the codomain by null functions
    (Subsection~\ref{sec:ae}).
    We do not, however, prove that ``it is undetermined unless one divides''
    (Problem~\ref{prob:no-lift}).
\end{itemize}

In the end we show that the Daniell extension gives a reflection in the category
with operator morphisms taken in this way as well
(Theorem~\ref{thm:operator-main-theorem}).

\subsection{Operator morphisms}

\begin{definition}\label{def:operator-morphism}
  Let $(X,L,I)$ and $(Y,M,J)$ be integration systems on sets $X$ and
  $Y$ respectively.
  An \textbf{operator morphism} of integration systems
  \[
    \Lambda:(X,L,I)\to(Y,M,J)
  \]
  is a linear map $\Lambda:L\to M$ satisfying the following two conditions.
  \begin{itemize}
    \item (positivity) for every $f\in L$, if $f\ge0$ then $\Lambda f\ge0$;
    \item (preservation of the integral) $J(\Lambda f)=I(f)$ for every $f\in L$.
  \end{itemize}
\end{definition}

\begin{remark}[relation to the operators of \cite{Yamagami}]\label{rem:yamagami-operator}
  \cite{Yamagami} introduces linear maps $\Lambda$ from a Riesz space $L$ on a
  set $X$ to a Riesz space $M$ on a set $Y$ satisfying
  \begin{itemize}
    \item (positivity) for every $f\in L$, if $f\ge0$ then $\Lambda f\ge0$;
    \item (order continuity) for every sequence of functions
      $\{f_n\}\subset L$, if $f_n(x)\downarrow0$ for all $x\in X$ then
      $(\Lambda f_n)(y)\downarrow0$ for all $y\in Y$.
  \end{itemize}
  Definition~\ref{def:operator-morphism} adds preservation of the integral to
  this and \textbf{removes order continuity}.

  Anything that is an operator in the sense of that book and preserves the
  integral is an operator morphism in the sense of
  Definition~\ref{def:operator-morphism}. The converse does not hold in general.
\end{remark}

\begin{remark}\label{rem:order-continuity-redundant}
  Note that we do not impose the order continuity of
  Remark~\ref{rem:yamagami-operator}.
  \cite{Yamagami} imposes it in order to \emph{construct} an integral
  $J\circ\Lambda$ on $L$ out of $J$ and $\Lambda$. Indeed, if
  $f_n(x)\downarrow0$ then $(\Lambda f_n)(y)\downarrow0$, and since $J$ is an
  integral, $J(\Lambda f_n)\to0$ follows.

  In this paper, by contrast, the integral $I$ is given from the outset as part
  of the object $(X,L,I)$.
  In that situation preservation of the integral, $J\circ\Lambda=I$,
  \textbf{already implies that $J\circ\Lambda$ is an integral}.
  That is, order continuity is unnecessary for the purpose of making
  $J\circ\Lambda$ an integral.

  Order continuity is a genuinely strong condition on the pointwise behaviour of
  $\Lambda f_n$, but it is never used in the discussion below.
  And in fact, for the morphisms of this category, pointwise conditions cannot
  be stated as propositions at all.
  The values of the extension given by
  Theorem~\ref{thm:operator-extension} are elements of the quotient
  $\widetilde{M}$, so a pointwise condition would depend on the choice of
  representative.
  That is, a pointwise condition is not false but meaningless.
  Imposing pointwise conditions at all is inappropriate in this category.

  Note also that the fact that the condition in \cite{Yamagami} is stated in
  terms of pointwise convergence is essential in that book's setting.
  If it were replaced by convergence in the order of $M$ then, since order
  convergence is strictly weaker than pointwise convergence (see the example in
  Appendix~\ref{sec:forget}), $J(\Lambda f_n)\to0$ would no longer follow.
\end{remark}

\begin{lemma}
  \begin{enumerate}
    \item For every integration system $(X,L,I)$, the identity map
      $\operatorname{id}_L:L\to L$ is an operator morphism
      $(X,L,I)\to(X,L,I)$.
    \item For operator morphisms $\Lambda:(X,L,I)\to(Y,M,J)$ and
      $\Xi:(Y,M,J)\to(Z,N,K)$, the composite $\Xi\circ\Lambda:L\to N$ is an
      operator morphism $(X,L,I)\to(Z,N,K)$.
  \end{enumerate}
\end{lemma}
\begin{proof}
  (1) is immediate. As for (2), a composite of linear maps is linear and a
  composite of positive maps is positive.
  Preservation of the integral holds because
  $K(\Xi\Lambda f)=J(\Lambda f)=I(f)$.
\end{proof}

\begin{definition}
  The category whose objects are integration systems and whose morphisms are
  operator morphisms of integration systems is written
  \[
    \mathbf{Int}_{\mathrm{op}} .
  \]
\end{definition}

\subsubsection{Relation to morphisms given by maps of base spaces}

\begin{lemma}\label{lem:base-map-gives-operator}
  Let $F:(X,L,I)\to(Y,M,J)$ be a morphism of integration systems in the sense of
  Section~\ref{chap:preliminaries}.
  Then
  \[
    F^*:L\to M;\quad f\mapsto f\circ\underline{F}
  \]
  is an operator morphism of integration systems.
\end{lemma}
\begin{proof}
  Linearity is immediate from the definition of composition.
  If $f\ge0$ then $(F^*f)(y)=f(\underline{F}(y))\ge0$ for every $y\in Y$, so
  $F^*$ is positive.
  Preservation of the integral is precisely the definition of a morphism of
  integration systems.
  Note that $F^*$ is order continuous as well. Indeed, if $f_n(x)\downarrow0$
  holds for every $x\in X$ then, putting $x:=\underline{F}(y)$ for each
  $y\in Y$, we get $(F^*f_n)(y)=f_n(x)\downarrow0$.
\end{proof}

\begin{proposition}\label{prop:base-to-operator-functor}
  The assignment that leaves objects unchanged and sends a morphism $F$ to
  $F^*$ defines a functor
  \[
    \mathcal{E}:\mathbf{Int}\to\mathbf{Int}_{\mathrm{op}} .
  \]
\end{proposition}
\begin{proof}
  By Lemma~\ref{lem:base-map-gives-operator} morphisms are sent to operator
  morphisms.
  For identity morphisms, $(\operatorname{id}_X)^*=\operatorname{id}_L$.
  As for composition, the map of base spaces underlying $G\circ F$ was
  $\underline{F}\circ\underline{G}$, so for every $f\in L$
  \[
    (G\circ F)^*f
    =f\circ(\underline{F}\circ\underline{G})
    =(f\circ\underline{F})\circ\underline{G}
    =G^*(F^*f) ,
  \]
  and $(G\circ F)^*=G^*\circ F^*$ holds.
\end{proof}

\begin{remark}
  The functor $\mathcal{E}$ need not be faithful.
  Indeed, when $L$ does not separate the points of $X$, distinct maps of base
  spaces $\underline{F}\ne\underline{G}$ can give the same pullback
  $F^*=G^*$.
  In this sense $\mathbf{Int}_{\mathrm{op}}$ does not carry from the outset the
  information about the base space that the integral does not distinguish.
\end{remark}

\begin{remark}
  $\mathcal{E}$ is not full either.
  Among operator morphisms, those coming from maps of base spaces are confined
  to the ones preserving the lattice operations $\wedge,\vee$ (Riesz
  homomorphisms).
  A general operator morphism need not be a lattice homomorphism, and, as
  \cite{Yamagami} points out, this difference is what decides whether
  probability kernels and Markov operators are included.
  The next example exhibits this difference concretely.
\end{remark}

\begin{example}[the case of finite sets: column-stochastic matrices]\label{ex:stochastic-matrix}
  Let $X,Y$ be finite sets and let $L=\mathbb{R}^X$ and $M=\mathbb{R}^Y$.
  Take positive weights $p:X\to(0,\infty)$ and $q:Y\to(0,\infty)$ and put
  \[
    I(f):=\sum_{x\in X}p(x)f(x),
    \qquad
    J(g):=\sum_{y\in Y}q(y)g(y) ;
  \]
  then $(X,L,I)$ and $(Y,M,J)$ are integration systems.

  A linear map $\Lambda:L\to M$ is represented by a matrix $(\Lambda_{yx})$.
  Then
  \begin{itemize}
    \item positivity is equivalent to all entries satisfying
      $\Lambda_{yx}\ge0$;
    \item preservation of the integral, $J(\Lambda f)=I(f)$, says that
      $\sum_{x}\bigl(\sum_{y}q(y)\Lambda_{yx}\bigr)f(x)=\sum_x p(x)f(x)$ holds
      for every $f$, hence is equivalent to
      \[
        \sum_{y\in Y}q(y)\Lambda_{yx}=p(x)
        \qquad(\forall x\in X) .
      \]
  \end{itemize}

  In particular, in the case $p\equiv1$ and $q\equiv1$ (counting, the finite
  version of Example~\ref{ex:counting}), the latter condition says
  \[
    \sum_{y\in Y}\Lambda_{yx}=1
    \qquad(\forall x\in X) ,
  \]
  that is, that \textbf{each column sums to $1$}.
  Hence in this case the operator morphisms of integration systems are precisely
  the \textbf{column-stochastic matrices}.
  Below, ``stochastic matrix'' in this paper means a column-stochastic matrix.

  On the other hand, for a morphism coming from a map of base spaces
  $\underline{F}:Y\to X$ we have $\Lambda_{yx}=1$ when
  $\underline{F}(y)=x$ and $\Lambda_{yx}=0$ otherwise, so each row contains
  exactly one $1$.
  In the counting case, moreover, the column-sum condition forces
  $|\underline{F}^{-1}(x)|=1$, that is, that $\underline{F}$ be a bijection.

  That is, between the same objects,
  \begin{center}
    with maps of base spaces as morphisms only the ``permutation matrices'',\\
    with operators as morphisms all ``column-stochastic matrices''
  \end{center}
  are morphisms. This is the simplest manifestation of the fact that
  $\mathcal{E}$ is not full.
\end{example}

\begin{remark}
  Imposing $\Lambda1_X=1_Y$ (row sums equal to $1$) in the setting of
  Example~\ref{ex:stochastic-matrix} gives what \cite{Yamagami} calls a
  \textbf{Markov operator}, the corresponding family of measures being a
  \textbf{transition probability}.

  Note, however, that in the counting case imposing preservation of the integral
  (column sums equal to $1$) and row sums equal to $1$ simultaneously yields a
  \textbf{doubly stochastic matrix}.
  Hence one cannot say that ``$\mathbf{Int}_{\mathrm{op}}$ contains all
  transition matrices of Markov chains''.
  For general weights $p,q$, the transition matrices that are morphisms are
  confined to those with $\sum_{y}q(y)\Lambda_{yx}=p(x)$, that is, those sending
  $q$ to $p$.
  What is included is exactly those satisfying this condition.
\end{remark}

\subsection{The Daniell extension of operator morphisms}

The main result of this subsection is
Theorem~\ref{thm:operator-extension}.
Unlike in Section~\ref{chap:functoriality},
\textbf{uniqueness does not follow from the agreement of maps of base spaces};
it is obtained by way of contractivity and the density of $L$ in $L^1$.

In what follows, for $(Y,M,J)$ the symbols $\mathcal{N}_J$,
$\widetilde{M}=M/\mathcal{N}_J$, $\pi_J$ and $\widetilde{J}$ are those of
Subsection~\ref{sec:ae}.

\subsubsection{Contractivity}

Below, we say that $\Lambda:L\to\widetilde{M}$ is an \textbf{operator morphism
into the quotient} if $\Lambda$ is linear, positive
($f\ge0\Rightarrow\Lambda f\ge0$) and satisfies
$\widetilde{J}(\Lambda f)=I(f)$.
For an ordinary operator morphism $\Lambda:L\to M$, the composite
$\pi_J\circ\Lambda$ is an operator morphism into the quotient.

\begin{lemma}\label{lem:operator-dominates-abs}
  Let $\Lambda:L\to\widetilde{M}$ be linear and positive.
  Then $|\Lambda f|\le\Lambda|f|$ holds for every $f\in L$.
\end{lemma}
\begin{proof}
  We have $f\le|f|$ and $-f\le|f|$. Since $\Lambda$ is linear and positive,
  $\Lambda f\le\Lambda|f|$ and $-\Lambda f\le\Lambda|f|$, and in the Riesz space
  $\widetilde{M}$ we obtain
  $|\Lambda f|=(\Lambda f)\vee(-\Lambda f)\le\Lambda|f|$.
\end{proof}

\begin{lemma}[contractivity]\label{lem:operator-is-contraction}
  Let $\Lambda:L\to\widetilde{M}$ be an operator morphism into the quotient.
  Then
  \[
    \|\Lambda f\|\le\|f\|_I
  \]
  holds for every $f\in L$. Moreover equality holds when $f\ge0$.
\end{lemma}
\begin{proof}
  By Lemma~\ref{lem:operator-dominates-abs} and positivity of $\widetilde{J}$,
  \[
    \|\Lambda f\|=\widetilde{J}(|\Lambda f|)\le\widetilde{J}(\Lambda|f|)=I(|f|)=\|f\|_I.
  \]
  Here the second equality is preservation of the integral for $|f|\in L$.
  When $f\ge0$ we have $|f|=f$ and $\Lambda f\ge0$, so
  $\|\Lambda f\|=\widetilde{J}(\Lambda f)=I(f)=\|f\|_I$.
\end{proof}

\begin{remark}
  That equality does not hold in general is the difference from a lattice
  homomorphism.
  If $\Lambda$ is a lattice homomorphism then $|\Lambda f|=\Lambda|f|$ and it is
  isometric, whereas a general operator morphism remains merely contractive.
  All that the discussion below needs is contractivity.
\end{remark}

\subsubsection{The extension theorem}

\begin{theorem}[Daniell extension of operator morphisms]\label{thm:operator-extension}
  Let $(X,L,I)$ be an integration system with Daniell extension
  $(X,L^1,I^1)$, and let $(Y,M,J)$ be an integration system
  \textbf{such that $\widetilde{M}$ is a Banach space}.
  Let $\Lambda:L\to\widetilde{M}$ be an operator morphism into the quotient.
  Then there \textbf{exists a unique} operator morphism into the quotient
  \[
    \Lambda^1:L^1\to\widetilde{M}
  \]
  satisfying $\Lambda^1|_L=\Lambda$.
  Moreover $\Lambda^1$ is contractive with respect to $\|\cdot\|$.
\end{theorem}
\begin{proof}
  \textbf{Existence.}
  Let $f\in L^1$. By Lemma~\ref{lem:l1-approx-by-source} we may take
  $\{s_N\}\subset L$ with $\|f-s_N\|_{I^1}\to0$.
  A convergent sequence is a Cauchy sequence, and
  $\|s_N-s_M\|_{I^1}=\|s_N-s_M\|_I$ for $s_N,s_M\in L$
  (Lemma~\ref{lem:daniell-ext-contains-source}).
  Hence Lemma~\ref{lem:operator-is-contraction} gives
  \[
    \|\Lambda s_N-\Lambda s_M\|
    =\|\Lambda(s_N-s_M)\|
    \le\|s_N-s_M\|_I
    \longrightarrow0 ,
  \]
  so $\{\Lambda s_N\}$ is a Cauchy sequence in $\widetilde{M}$.
  By hypothesis $\widetilde{M}$ is a Banach space, so this sequence converges.
  We define $\Lambda^1(f)$ to be its limit.

  \textbf{Independence of the approximating sequence.}
  If $\{t_N\}\subset L$ also satisfies $\|f-t_N\|_{I^1}\to0$ then
  $\|s_N-t_N\|_{I^1}\to0$, and contractivity gives
  $\|\Lambda s_N-\Lambda t_N\|\to0$.
  As $\widetilde{M}$ is a normed space and limits are unique, the two limits
  agree.

  \textbf{Linearity, contractivity, and agreement on $L$.}
  For $f,g\in L^1$ and $a\in\mathbb{R}$, taking approximating sequences
  $\{s_N\},\{u_N\}$ for each, the sequence $\{s_N+au_N\}$ is an approximating
  sequence for $f+ag$, whence linearity.
  Moreover $\|\Lambda s_N\|\le\|s_N\|_I=\|s_N\|_{I^1}\to\|f\|_{I^1}$, so
  continuity of the norm gives $\|\Lambda^1 f\|\le\|f\|_{I^1}$.
  For $f\in L$, taking the constant sequence $s_N:=f$ gives
  $\Lambda^1 f=\Lambda f$.

  \textbf{Positivity.}
  Let $f\in L^1$ with $f\ge0$.
  For an approximating sequence $\{s_N\}$ put $s_N^+:=s_N\vee0\in L$.
  Since $f\ge0$ we have $|f-s_N^+|\le|f-s_N|$ pointwise, hence
  $\|f-s_N^+\|_{I^1}\le\|f-s_N\|_{I^1}\to0$.
  So $\{s_N^+\}$ is also an approximating sequence for $f$, and
  $\Lambda^1 f=\lim_N\Lambda s_N^+$.
  For each $N$, $s_N^+\ge0$ gives $\Lambda s_N^+\ge0$.
  The positive cone of $\widetilde{M}$ is closed with respect to $\|\cdot\|$
  (indeed $\|(\Lambda^1 f)\wedge0\|\le\|\Lambda^1 f-\Lambda s_N^+\|\to0$), so
  $\Lambda^1 f\ge0$.

  \textbf{Preservation of the integral.}
  $\widetilde{J}$ satisfies $|\widetilde{J}(v)|\le\|v\|$ and is therefore
  continuous.
  Hence $\widetilde{J}(\Lambda^1 f)=\lim_N\widetilde{J}(\Lambda s_N)
  =\lim_N I(s_N)=\lim_N I^1(s_N)=I^1(f)$.
  The last equality holds because
  $|I^1(f)-I^1(s_N)|\le I^1(|f-s_N|)\to0$.

  \textbf{Uniqueness.}
  Suppose $\Psi:L^1\to\widetilde{M}$ is an operator morphism into the quotient
  with $\Psi|_L=\Lambda$.
  Applying Lemma~\ref{lem:operator-is-contraction} to $(X,L^1,I^1)$, the map
  $\Psi$ is contractive and hence continuous with respect to
  $\|\cdot\|_{I^1}$.
  $\Psi$ and $\Lambda^1$ agree on $L$, and $L$ is dense in $L^1$ by
  Lemma~\ref{lem:l1-approx-by-source}, so as two continuous maps they agree on
  all of $L^1$.
\end{proof}

\begin{remark}\label{rem:uniqueness-has-content}
  Note that the proof of uniqueness differs essentially from the case of the
  main theorem of Section~\ref{chap:functoriality}.
  There, morphisms were the maps of base spaces themselves, so uniqueness of the
  factorizing morphism followed automatically from the agreement of the maps of
  base spaces.
  Here \textbf{contractivity
  (Lemma~\ref{lem:operator-is-contraction}) and density
  (Lemma~\ref{lem:l1-approx-by-source}) actually do the work.}

  Section~\ref{chap:operator} is not, however, the first place where an analytic
  lemma from the preliminary section contributes to a proof of a universal
  property.
  In the ae version of the main theorem in
  Subsection~\ref{sec:main-theorem-ae}, the same density lemma is used by way of
  Lemma~\ref{lem:daniell-extension-preserves-ae}.
  Precisely, \textbf{in the non-ae version uniqueness follows automatically from
  the agreement of the maps of base spaces, while the ae version and the
  operator version of Section~\ref{chap:operator} require an analytic
  argument}.
\end{remark}

\begin{remark}\label{rem:ae-is-forced}
  What Theorem~\ref{thm:operator-extension} gives is a map into
  $\widetilde{M}$, not a map into $M$.
  Choosing representatives in $M$ for the elements of $\widetilde{M}$ yields a
  map $L^1\to M$, but this is determined only up to null functions.
  Hence, in order to discuss the Daniell extension in a category whose morphisms
  are operator morphisms, it is natural to divide the codomain by null
  functions.

  This paper cannot, however, assert that ``the extension is undetermined unless
  one divides''.
  That would require an \textbf{example of an operator morphism into the
  quotient admitting no linear positive lift $L^1\to M$}, and this paper does
  not supply one (Problem~\ref{prob:no-lift}).
  Indeed, for morphisms coming from maps of base spaces a lift
  $F^*:L^1\to M$ does exist, so there is no need to pass to ae-equivalence in
  the setting of Section~\ref{chap:functoriality}.

  Note that the situation in which the extension is determined only up to null
  sets occurs likewise in the theorem on iterated integrals in
  \cite{Yamagami}, and that book states the extension for integrable functions
  in that form.
\end{remark}

\begin{problem}\label{prob:no-lift}
  Is there an operator morphism into the quotient $L^1\to\widetilde{M}$
  admitting no linear positive lift $L^1\to M$?

  If one exists, it would show that passing to ae-equivalence is a requirement,
  and at the same time add one more reason to keep
  Section~\ref{chap:functoriality} as an independent section (see
  Subsection~\ref{sec:two-chapters}).
\end{problem}

\subsubsection{The ae-category and the main theorem}

\begin{lemma}[an operator morphism into the quotient is determined by the quotient]\label{lem:operator-descends}
  Let $\Lambda:L\to\widetilde{M}$ be an operator morphism into the quotient.
  Then $\Lambda$ maps $\mathcal{N}_I$ to $0$, and there is a unique linear map
  $\widetilde{\Lambda}:\widetilde{L}\to\widetilde{M}$ satisfying
  $\widetilde{\Lambda}\circ\pi_I=\Lambda$.
  Moreover $\widetilde{\Lambda}$ is positive and satisfies
  $\widetilde{J}(\widetilde{\Lambda}(\pi_I(f)))=\widetilde{I}(\pi_I(f))$ for
  every $f\in L$.
\end{lemma}
\begin{proof}
  If $f\in\mathcal{N}_I$ then $\|f\|_I=0$, and
  Lemma~\ref{lem:operator-is-contraction} gives
  $\|\Lambda f\|\le\|f\|_I=0$.
  As $\|\cdot\|$ is a norm on $\widetilde{M}$ we get $\Lambda f=0$.
  Hence $\Lambda$ maps $\mathcal{N}_I$ to $0$ and a unique linear map
  $\widetilde{\Lambda}$ is determined.

  We check positivity. If $\pi_I(f)\ge0$ then, by the way the order on
  $\widetilde{L}$ is defined (Lemma~\ref{lem:quotient-riesz}), we may take
  $f'\in\mathcal{N}_I$ with $f+f'\ge0$.
  Since $\Lambda f'=0$,
  \[
    \widetilde{\Lambda}(\pi_I(f))=\Lambda(f+f')\ge0 .
  \]
  Finally
  $\widetilde{J}(\widetilde{\Lambda}(\pi_I(f)))=\widetilde{J}(\Lambda f)=I(f)=\widetilde{I}(\pi_I(f))$.
\end{proof}

\begin{definition}\label{def:int-op-ae}
  The category whose objects are integration systems and whose morphisms from
  $(X,L,I)$ to $(Y,M,J)$ are the operator morphisms into the quotient
  $\Lambda:L\to\widetilde{M}$ is written
  \[
    \mathbf{Int}_{\mathrm{op},\mathrm{ae}} .
  \]
  Composition is defined as follows.
  For $\Lambda:L\to\widetilde{M}$ and $\Xi:M\to\widetilde{N}$,
  Lemma~\ref{lem:operator-descends} shows that $\Xi$ induces a linear positive
  map $\widetilde{\Xi}:\widetilde{M}\to\widetilde{N}$.
  We define the composite to be $\widetilde{\Xi}\circ\Lambda$.
  This is linear and positive and satisfies
  $\widetilde{K}(\widetilde{\Xi}(\Lambda f))=\widetilde{J}(\Lambda f)=I(f)$, so
  it is again an operator morphism into the quotient.
  The identity morphism is $\pi_I:L\to\widetilde{L}$.

  We write the full subcategory of objects for which $\widetilde{L}$ is a Banach
  space as
  \[
    \mathbf{BInt}_{\mathrm{op},\mathrm{ae}} .
  \]
  Daniell closed integration systems belong to this subcategory by
  Lemma~\ref{lem:quotient-is-banach}.
  The reason for taking this subcategory is given in
  Subsection~\ref{sec:what-op-ae-sees}.
\end{definition}

\begin{remark}
  For ordinary operator morphisms $\Lambda,\Lambda':L\to M$, defining the same
  morphism $\pi_J\circ\Lambda=\pi_J\circ\Lambda'$ is equivalent to
  $J(|\Lambda f-\Lambda' f|)=0$ holding for every $f\in L$.
  The latter is precisely the relation $=_{\mathrm{ae}}$ of
  Definition~\ref{def:ae-equal} transported to operator morphisms.

  It is not accurate, however, to call
  $\mathbf{Int}_{\mathrm{op},\mathrm{ae}}$ ``the category obtained by dividing
  morphisms by ae-equivalence''.
  Its morphisms are defined from the outset as \emph{all} operator morphisms
  into the quotient $L\to\widetilde{M}$, and not all of them need come from
  $M$-valued morphisms (Remark~\ref{rem:ae-is-forced}).
  Precisely, the relation is that \textbf{there is a faithful functor from the
  category obtained by dividing ordinary operator morphisms by ae-equivalence
  into $\mathbf{Int}_{\mathrm{op},\mathrm{ae}}$}.
\end{remark}

\begin{theorem}[operator-morphism version of the main theorem]\label{thm:operator-main-theorem}
  The inclusion functor
  \[
    \iota_{B,\mathrm{op}}:
    \mathbf{BInt}_{\mathrm{op},\mathrm{ae}}
    \to
    \mathbf{Int}_{\mathrm{op},\mathrm{ae}}
  \]
  has a left adjoint
  \[
    \mathcal{D}_{\mathrm{op}}:
    \mathbf{Int}_{\mathrm{op},\mathrm{ae}}
    \to
    \mathbf{BInt}_{\mathrm{op},\mathrm{ae}} .
  \]
  Its action on objects is the Daniell extension.
  Its action on morphisms sends a morphism $\Lambda:A\to B$ to the extension to
  $L^1\to\widetilde{M^1}$, given by
  Theorem~\ref{thm:operator-extension}, of the composite
  $i_B\circ\Lambda:L\to\widetilde{M^1}$
  (since $\mathcal{D}B$ is Daniell closed, $\widetilde{M^1}$ is a Banach space
  by Lemma~\ref{lem:quotient-is-banach} and that theorem applies).
\end{theorem}
\begin{proof}
  Let $A=(X,L,I)$ be an object of
  $\mathbf{Int}_{\mathrm{op},\mathrm{ae}}$ and let
  $\mathcal{D}A:=(X,L^1,I^1)$ be its Daniell extension.
  By Lemma~\ref{lem:daniell-ext-is-complete} the system $\mathcal{D}A$ is
  Daniell closed, and by Lemma~\ref{lem:quotient-is-banach} the space
  $\widetilde{L^1}$ is a Banach space, so $\mathcal{D}A$ is an object of
  $\mathbf{BInt}_{\mathrm{op},\mathrm{ae}}$.
  Moreover, from the inclusion $L\subset L^1$ and $I^1|_L=I$
  (Lemma~\ref{lem:daniell-ext-contains-source}), the map
  \[
    i_A:=\pi_{I^1}|_L:L\to\widetilde{L^1}
  \]
  is an operator morphism into the quotient and defines a morphism
  $i_A:A\to\iota_{B,\mathrm{op}}(\mathcal{D}A)$ of
  $\mathbf{Int}_{\mathrm{op},\mathrm{ae}}$.

  We show that $i_A$ is a universal morphism from $A$ to
  $\iota_{B,\mathrm{op}}$.
  Let $B=(Y,M,J)$ be an object of
  $\mathbf{BInt}_{\mathrm{op},\mathrm{ae}}$ and let
  $\Lambda:A\to\iota_{B,\mathrm{op}}(B)$ be a morphism, that is, an operator
  morphism into the quotient $\Lambda:L\to\widetilde{M}$.
  Since $\widetilde{M}$ is a Banach space,
  Theorem~\ref{thm:operator-extension} applies and there is a unique operator
  morphism into the quotient $\Lambda^1:L^1\to\widetilde{M}$ with
  $\Lambda^1|_L=\Lambda$.

  Write $\Lambda^\sharp$ for the morphism $\mathcal{D}A\to B$ determined by
  $\Lambda^1$.
  By the definition of composition in
  Definition~\ref{def:int-op-ae}, the composite
  $\iota_{B,\mathrm{op}}(\Lambda^\sharp)\circ i_A$ sends $f\in L$ to
  $\widetilde{\Lambda^1}(\pi_{I^1}(f))=\Lambda^1(f)=\Lambda(f)$.
  Hence
  \[
    \Lambda=\iota_{B,\mathrm{op}}(\Lambda^\sharp)\circ i_A.
  \]

  Conversely, if a morphism $H:\mathcal{D}A\to B$ satisfies
  $\Lambda=\iota_{B,\mathrm{op}}(H)\circ i_A$, then the operator morphism into
  the quotient $L^1\to\widetilde{M}$ representing $H$ agrees with $\Lambda$ on
  $L$.
  By the uniqueness in Theorem~\ref{thm:operator-extension} we get
  $H=\Lambda^\sharp$.

  Since $A$ was arbitrary,
  Theorem~\ref{thm:left-adjunction-existence-from-universality} shows that
  $\iota_{B,\mathrm{op}}$ has a left adjoint.
\end{proof}

\begin{remark}
  What Section~\ref{chap:operator} imposes on the codomain is only that
  $\widetilde{M}$ be a Banach space, not Daniell closedness itself.
  For Daniell closed integration systems,
  Lemma~\ref{lem:quotient-is-banach} bridges the two.
  That this choice of subcategory is the natural one is seen in the next
  subsubsection (Proposition~\ref{prop:BInt-characterization}).

  Note also that the definition of the monotone convergence theorem
  (Definition~\ref{def:satisfies-mct}) is used with pointwise convergence
  unchanged.
  That is a property of \emph{objects} only, and it works solely to yield
  completeness of $\widetilde{M}$ by way of
  Lemma~\ref{lem:quotient-is-banach}.
  As long as no pointwise condition is imposed on the side of morphisms
  (Remark~\ref{rem:order-continuity-redundant}), there is no need to rewrite the
  monotone convergence theorem in an almost-everywhere sense.
\end{remark}

\subsubsection{What this category sees}\label{sec:what-op-ae-sees}

There is a point about the formulation of the main theorem of the previous
subsubsection that deserves attention.
The morphisms of $\mathbf{Int}_{\mathrm{op},\mathrm{ae}}$ are operator morphisms
into the quotient $L\to\widetilde{M}$, and by contractivity
(Lemma~\ref{lem:operator-is-contraction}) they map $\mathcal{N}_I$ to $0$.
Hence morphisms see only $\widetilde{L}$ and $\widetilde{M}$.
In this subsubsection we formulate this and state what follows from it.

\begin{definition}\label{def:preAL}
  We say that a pair $(E,\phi)$ is a \textbf{pre-AL space} if $E$ is a Riesz
  space and $\phi:E\to\mathbb{R}$ is a positive linear functional such that
  \[
    \|u\|_\phi:=\phi(|u|)
  \]
  is a norm on $E$.
  A morphism from $(E,\phi)$ to $(F,\psi)$ is a linear map $T:E\to F$ that is
  positive and satisfies $\psi\circ T=\phi$.
  We write this category as $\mathbf{AL}_0$.
\end{definition}

\noindent
Since $\|u+v\|_\phi=\phi(u+v)=\|u\|_\phi+\|v\|_\phi$ for $u,v\ge0$, the norm is
additive on the positive cone.
If it is complete this is an AL-space, the object of Kakutani's representation
theorem touched on in Appendix~\ref{sec:forget}.

\begin{lemma}\label{lem:quotient-is-preAL}
  If $(X,L,I)$ is an integration system then
  $(\widetilde{L},\widetilde{I})$ is a pre-AL space.
\end{lemma}
\begin{proof}
  By Lemma~\ref{lem:quotient-riesz}.
\end{proof}

\begin{proposition}\label{prop:Q-full-faithful}
  The assignment sending an object $A=(X,L,I)$ to
  $(\widetilde{L},\widetilde{I})$ and a morphism $\Lambda:L\to\widetilde{M}$ to
  the induced map $\widetilde{\Lambda}:\widetilde{L}\to\widetilde{M}$ defines a
  \textbf{faithful and full} functor
  \[
    Q:\mathbf{Int}_{\mathrm{op},\mathrm{ae}}\longrightarrow\mathbf{AL}_0 .
  \]
\end{proposition}
\begin{proof}
  \textbf{That $\widetilde{\Lambda}$ is a morphism of $\mathbf{AL}_0$} is
  Lemma~\ref{lem:operator-descends}.

  \textbf{Functoriality.}
  The identity morphism is $\pi_I$ and
  $\widetilde{\pi_I}=\operatorname{id}_{\widetilde{L}}$.
  As for composition, Definition~\ref{def:int-op-ae} gives
  $\Xi\circ\Lambda=\widetilde{\Xi}\circ\Lambda
  =\widetilde{\Xi}\circ\widetilde{\Lambda}\circ\pi_I$, so
  $\widetilde{\Xi\circ\Lambda}=\widetilde{\Xi}\circ\widetilde{\Lambda}$.

  \textbf{Faithfulness.}
  If $\widetilde{\Lambda}=\widetilde{\Lambda'}$ then
  $\Lambda=\widetilde{\Lambda}\circ\pi_I=\widetilde{\Lambda'}\circ\pi_I=\Lambda'$.

  \textbf{Fullness.}
  Given a morphism
  $T:(\widetilde{L},\widetilde{I})\to(\widetilde{M},\widetilde{J})$ of
  $\mathbf{AL}_0$, the map $\Lambda:=T\circ\pi_I:L\to\widetilde{M}$ is linear
  and positive and satisfies
  $\widetilde{J}(\Lambda f)=\widetilde{I}(\pi_I(f))=I(f)$, so it is an operator
  morphism into the quotient.
  By construction $\widetilde{\Lambda}=T$.
\end{proof}

\begin{corollary}\label{cor:op-ae-iso-criterion}
  $A\cong B$ in $\mathbf{Int}_{\mathrm{op},\mathrm{ae}}$ is equivalent to
  $(\widetilde{L},\widetilde{I})\cong(\widetilde{M},\widetilde{J})$ in
  $\mathbf{AL}_0$.
  Moreover the isomorphism is then isometric.
\end{corollary}
\begin{proof}
  A faithful and full functor reflects isomorphisms, and also preserves them.
  As for isometry, since both $T$ and its inverse $T^{-1}$ are contractive (use
  Lemma~\ref{lem:operator-is-contraction} with the same computation for
  morphisms of $\mathbf{AL}_0$),
  \[
    \|u\|=\|T^{-1}Tu\|\le\|Tu\|\le\|u\|
  \]
  follows.
\end{proof}

\noindent
That is, \textbf{$\mathbf{Int}_{\mathrm{op},\mathrm{ae}}$ sees neither the base
space $X$ nor the realization of $\widetilde{L}$ as a function space.}
If $\widetilde{L}$ is the same as a pre-AL space, the objects are not
distinguished.

\begin{proposition}\label{prop:D-op-is-banach-completion}
  Let $A=(X,L,I)$ be an integration system. Then
  $Q(i_A)=j_I:\widetilde{L}\to\widetilde{L^1}$ is a dense isometric embedding
  and $\widetilde{L^1}$ is a Banach space.
  That is, $Q(\mathcal{D}A)$ is the Banach completion of $Q(A)$.
\end{proposition}
\begin{proof}
  That $j_I$ is an injective Riesz homomorphism and is isometric is
  Lemma~\ref{lem:quotient-embedding}.
  Density of the image: for $f\in L^1$, taking $\{s_N\}\subset L$ as in
  Lemma~\ref{lem:l1-approx-by-source},
  \[
    \|\pi_{I^1}(f)-j_I(\pi_I(s_N))\|
    =\|\pi_{I^1}(f-s_N)\|
    =I^1(|f-s_N|)
    \longrightarrow0.
  \]
  Since $\mathcal{D}A$ is a Daniell closed integration system
  (Lemma~\ref{lem:daniell-ext-is-complete}),
  Lemma~\ref{lem:quotient-is-banach} shows that $\widetilde{L^1}$ is a Banach
  space.
\end{proof}

\begin{proposition}[characterization of $\mathbf{BInt}_{\mathrm{op},\mathrm{ae}}$]\label{prop:BInt-characterization}
  Let $A=(X,L,I)$ be an object of
  $\mathbf{Int}_{\mathrm{op},\mathrm{ae}}$.
  Then $A$ belonging to $\mathbf{BInt}_{\mathrm{op},\mathrm{ae}}$ is equivalent
  to $i_A:A\to\mathcal{D}A$ being an isomorphism of
  $\mathbf{Int}_{\mathrm{op},\mathrm{ae}}$.
  Moreover $\mathbf{BInt}_{\mathrm{op},\mathrm{ae}}$ is isomorphism-closed, and
  is the isomorphism closure of the class of Daniell closed objects.
\end{proposition}
\begin{proof}
  $\Leftarrow$: if $i_A$ is an isomorphism then
  Corollary~\ref{cor:op-ae-iso-criterion} shows that $\widetilde{L}$ is
  isometrically isomorphic to $\widetilde{L^1}$, and the latter is a Banach
  space by Proposition~\ref{prop:D-op-is-banach-completion}.

  $\Rightarrow$: if $\widetilde{L}$ is a Banach space then the image of the
  isometric embedding $j_I$ of
  Proposition~\ref{prop:D-op-is-banach-completion} is complete, hence closed,
  and it is also dense, so $j_I$ is surjective.
  As $j_I$ is a Riesz isomorphism, its inverse is positive and preserves the
  integral, so $Q(i_A)$ is an isomorphism of $\mathbf{AL}_0$.
  By Proposition~\ref{prop:Q-full-faithful} the morphism $i_A$ is an
  isomorphism.

  Being isomorphism-closed follows from
  Corollary~\ref{cor:op-ae-iso-criterion}.
  Isomorphisms are isometric, so if $\widetilde{L}$ is a Banach space then any
  pre-AL space isomorphic to it is a Banach space too.

  Finally we check that it is the closure.
  Daniell closed objects belong to
  $\mathbf{BInt}_{\mathrm{op},\mathrm{ae}}$ by
  Lemma~\ref{lem:quotient-is-banach}, and as just seen that subcategory is
  isomorphism-closed, so the isomorphism closure of the Daniell closed objects
  is contained in $\mathbf{BInt}_{\mathrm{op},\mathrm{ae}}$.
  Conversely, if $A\in\mathbf{BInt}_{\mathrm{op},\mathrm{ae}}$ then
  $A\cong\mathcal{D}A$, and $\mathcal{D}A$ is Daniell closed.
\end{proof}

\begin{remark}[the class of Daniell closed objects is not isomorphism-closed]\label{rem:op-ae-not-replete}
  The proposition above means that the subcategory to be taken when stating the
  main theorem in this category is
  $\mathbf{BInt}_{\mathrm{op},\mathrm{ae}}$, not the class of all Daniell closed
  objects.
  The latter is not isomorphism-closed.

  Indeed, take the example of
  Remark~\ref{rem:D-I-not-left-adjoint}: $X=\{0,1,2\}$,
  $L=\{(a,b,b)\}$, $I(a,b,b)=a$.
  Since $I(|(a,b,b)|)=|a|$ we have
  $\mathcal{N}_I=\{(0,b,b)\mid b\in\mathbb{R}\}$ and
  \[
    \widetilde{L}\cong\mathbb{R},
    \qquad
    \widetilde{I}(t)=t
  \]
  (the isomorphism being $\pi_I(a,b,b)\mapsto a$).
  As $\mathbb{R}$ is a Banach space, the proposition above gives
  $A\cong\mathcal{D}A$.
  On the other hand, as seen in that example, $L^1=\mathbb{R}^X\ne L$, so $A$ is
  not Daniell closed.

  This is not a counterexample to the main theorem
  (Theorem~\ref{thm:operator-main-theorem}), since being a reflective
  subcategory does not require being isomorphism-closed.
  The point is that $\mathbf{BInt}_{\mathrm{op},\mathrm{ae}}$ is the more
  natural choice for this category.
\end{remark}

\begin{remark}[the position of this section]\label{rem:op-ae-is-banach-completion}
  To summarize, what the Daniell extension gives in
  $\mathbf{Int}_{\mathrm{op},\mathrm{ae}}$ is the \textbf{Banach completion of
  the pre-AL space $(\widetilde{L},\widetilde{I})$}
  (Proposition~\ref{prop:D-op-is-banach-completion}).
  This category sees neither the base space nor the realization as a function
  space, so \textbf{one cannot expect the main theorem of
  Section~\ref{chap:operator} to carry the content ``accompanied by a concrete
  realization''.}
  The distinction from an abstract completion described in
  Subsection~\ref{sec:prior-work} is lost in the category of
  Section~\ref{chap:operator}.

  That distinction is retained on the side of the category of
  Section~\ref{chap:functoriality}.
  There, morphisms are the maps of base spaces themselves, and objects too are
  distinguished together with their realization $L\subset\mathbb{R}^X$.
  For the relation between the two sections see
  Subsection~\ref{sec:two-chapters}.

  Section~\ref{chap:operator} does not thereby become devoid of content.
  That a Banach completion \emph{exists} is general theory, but what
  Section~\ref{chap:operator} states is that \textbf{that completion is given as
  $\widetilde{L^1}$, the quotient of the integrable functions on the same base
  space} (Proposition~\ref{prop:D-op-is-banach-completion}).
  This is an assertion at the level of objects and is invisible to the morphisms
  of $\mathbf{Int}_{\mathrm{op},\mathrm{ae}}$.
  In other words, the concrete realization survives \emph{within the statement
  of the theorem} but not \emph{as categorical structure}.

  There is one more reason to keep Section~\ref{chap:operator}.
  Operator morphisms include, as morphisms and without modification, operations
  occurring in probability theory.
  Conditional expectation $E[\cdot\mid\mathcal{G}]$ is linear and positive and
  preserves the integral, so it is exactly an operator morphism of
  Section~\ref{chap:operator} (and it satisfies order continuity in the sense of
  Remark~\ref{rem:yamagami-operator} as well).
  In the framework of Section~\ref{chap:functoriality}, whose morphisms are maps
  of base spaces, this operation cannot be expressed as a morphism.
  The entrance to treating the It\^o integral and martingales in the same
  framework (Appendix~\ref{sec:forget}) is on the side of
  Section~\ref{chap:operator}.
\end{remark}

\subsection{Relation to the case of maps of base spaces as morphisms}

Section~\ref{chap:functoriality} and Section~\ref{chap:operator} put different
morphisms on the same objects and obtained a reflection theorem of the same
shape for each.
In this subsection we clarify the relation between the two.

We state the conclusion first.
The functor $\mathcal{E}$ commutes with the two reflections
(Proposition~\ref{prop:E-commutes-with-D}), and as a consequence
\textbf{the uniqueness part of the main theorem of
Section~\ref{chap:functoriality} follows from the theorem of
Section~\ref{chap:operator}}
(Corollary~\ref{cor:ch3-uniqueness-from-ch4}).
The existence part, however, \textbf{does not follow}.
That is, one cannot say from the operator version alone that the extension comes
from a map of base spaces
(Remark~\ref{rem:two-chapters-relation}).

\subsubsection{A functor between the ae-categories}

\begin{lemma}\label{lem:E-descends-to-ae}
  For a morphism of integration systems
  $F:(X,L,I)\to(Y,M,J)$ put
  \[
    \mathcal{E}_{\mathrm{ae}}(F):=\pi_J\circ F^*:L\to\widetilde{M} .
  \]
  Then the following hold.
  \begin{enumerate}
    \item $\mathcal{E}_{\mathrm{ae}}(F)$ is an operator morphism into the
      quotient.
    \item $F=_{\mathrm{ae}}G$ and
      $\mathcal{E}_{\mathrm{ae}}(F)=\mathcal{E}_{\mathrm{ae}}(G)$ are
      equivalent.
    \item The assignment $[F]\mapsto\mathcal{E}_{\mathrm{ae}}(F)$ defines a
      \textbf{faithful} functor
      \[
        \mathcal{E}_{\mathrm{ae}}:
        \mathbf{Int}_{\mathrm{ae}}\to\mathbf{Int}_{\mathrm{op},\mathrm{ae}} ,
      \]
      which is the identity on objects.
  \end{enumerate}
\end{lemma}
\begin{proof}
  (1) By Lemma~\ref{lem:base-map-gives-operator} the map $F^*$ is linear and
  positive and preserves the integral, and $\pi_J$ is a linear Riesz
  homomorphism (Lemma~\ref{lem:quotient-riesz}), so the composite is linear and
  positive as well.
  Moreover $\widetilde{J}(\pi_J(F^*f))=J(F^*f)=I(f)$.

  (2) $\pi_J(F^*f)=\pi_J(G^*f)$ is equivalent to
  $F^*f-G^*f\in\mathcal{N}_J$, that is, to
  $J(|F^*f-G^*f|)=0$.
  That this holds for every $f\in L$ was the definition of
  $F=_{\mathrm{ae}}G$ (Definition~\ref{def:ae-equal}).

  (3) The ``only if'' direction of (2) shows that the assignment is independent
  of the choice of representative.
  The converse direction gives injectivity, that is, faithfulness.
  For identity morphisms,
  $\mathcal{E}_{\mathrm{ae}}(\operatorname{id})=\pi_I$, which is the identity
  morphism of $\mathbf{Int}_{\mathrm{op},\mathrm{ae}}$
  (Definition~\ref{def:int-op-ae}).
  As for composition, let $F:(X,L,I)\to(Y,M,J)$ and
  $G:(Y,M,J)\to(Z,N,K)$. By the definition of composition in
  Definition~\ref{def:int-op-ae}, the map
  $\widetilde{M}\to\widetilde{N}$ induced by
  $\mathcal{E}_{\mathrm{ae}}(G)$ is $\pi_J(h)\mapsto\pi_K(G^*h)$. Hence, for
  every $f\in L$,
  \[
    \bigl(\widetilde{\mathcal{E}_{\mathrm{ae}}(G)}\circ\mathcal{E}_{\mathrm{ae}}(F)\bigr)(f)
    =\pi_K\bigl(G^*(F^*f)\bigr)
    =\pi_K\bigl((G\circ F)^*f\bigr)
    =\mathcal{E}_{\mathrm{ae}}(G\circ F)(f) .
  \]
\end{proof}

\begin{remark}\label{rem:E-ae-faithful-not-full}
  As stated at the beginning of Section~\ref{chap:operator}, the functor
  $\mathcal{E}:\mathbf{Int}\to\mathbf{Int}_{\mathrm{op}}$ before passing to the
  ae-quotient need not be faithful, because distinct maps of base spaces can
  give the same pullback.
  Once the ae-quotient is taken, exactly that difference is collapsed by
  $=_{\mathrm{ae}}$, so $\mathcal{E}_{\mathrm{ae}}$ becomes faithful.

  On the other hand $\mathcal{E}_{\mathrm{ae}}$ is \textbf{not full}.
  Only the morphisms coming from lattice homomorphisms, that is, the
  deterministic ones, lie in the image, and a general operator morphism is not
  among them.
  In the case of finite sets, the image consists only of those
  column-stochastic matrices each of whose rows has exactly one $1$.
\end{remark}

\subsubsection{Commutativity of the two reflections}

We write $\mathcal{D}_{\mathrm{op}}$ for the left adjoint of
$\iota_{B,\mathrm{op}}$ given by
Theorem~\ref{thm:operator-main-theorem}.
Its action is as stated in Theorem~\ref{thm:operator-main-theorem}.

\begin{proposition}\label{prop:E-commutes-with-D}
  The following diagram commutes.
  \[
    \begin{array}{ccc}
      \mathbf{Int}_{\mathrm{ae}}
        & \xrightarrow{\ \mathcal{D}_{\mathrm{ae}}\ }
        & \mathbf{DInt}_{\mathrm{ae}}\\[2pt]
      \big\downarrow{\scriptstyle\mathcal{E}_{\mathrm{ae}}}
        & & \big\downarrow{\scriptstyle\mathcal{E}_{\mathrm{ae}}}\\[2pt]
      \mathbf{Int}_{\mathrm{op},\mathrm{ae}}
        & \xrightarrow{\ \mathcal{D}_{\mathrm{op}}\ }
        & \mathbf{BInt}_{\mathrm{op},\mathrm{ae}}
    \end{array}
  \]
  That is,
  $\mathcal{E}_{\mathrm{ae}}\circ\mathcal{D}_{\mathrm{ae}}
  =\mathcal{D}_{\mathrm{op}}\circ\mathcal{E}_{\mathrm{ae}}$.
\end{proposition}
\begin{proof}
  On objects both are the Daniell extension
  $(X,L,I)\mapsto(X,L^1,I^1)$, and $\mathcal{E}_{\mathrm{ae}}$ does not change
  objects, so they agree.

  We argue for morphisms. Let $[F]:(X,L,I)\to(X',L',I')$ be a morphism of
  $\mathbf{Int}_{\mathrm{ae}}$.
  Write the injective Riesz homomorphism determined by
  Lemma~\ref{lem:quotient-embedding} as
  \[
    j':\widetilde{L'}\longrightarrow\widetilde{L'^1} .
  \]
  The unit morphism
  $i_{A'}:(X',L',I')\to(X',L'^1,I'^1)$ in
  $\mathbf{Int}_{\mathrm{op},\mathrm{ae}}$ is
  $\pi_{I'^1}|_{L'}=j'\circ\pi_{I'}$.

  Since $\mathcal{D}_{\mathrm{op}}$ is a left adjoint, for an operator morphism
  into the quotient $\Lambda:L\to\widetilde{L'}$ the morphism
  $\mathcal{D}_{\mathrm{op}}(\Lambda)$ is the unique operator morphism into the
  quotient $L^1\to\widetilde{L'^1}$ satisfying
  \[
    \mathcal{D}_{\mathrm{op}}(\Lambda)\bigr|_{L}=j'\circ\Lambda .
  \]
  Applying this to
  $\Lambda=\mathcal{E}_{\mathrm{ae}}[F]=\pi_{I'}\circ F^*$ gives the composite
  going down then across in the diagram, which we write as $\Phi_2$.

  The composite going across then down, on the other hand, is
  $\mathcal{E}_{\mathrm{ae}}$ applied to
  $\mathcal{D}_{\mathrm{ae}}[F]=[F^1]$, that is,
  \[
    \Phi_1:=\pi_{I'^1}\circ(F^1)^*:L^1\to\widetilde{L'^1} .
  \]
  Applying Lemma~\ref{lem:E-descends-to-ae} to the morphism $F^1$, the map
  $\Phi_1$ is an operator morphism into the quotient.

  We check that $\Phi_1$ satisfies the same characterization as $\Phi_2$.
  For $f\in L$ we have $(F^1)^*f=F^*f\in L'$, so the equality
  $\pi_{I'^1}|_{L'}=j'\circ\pi_{I'}$ of
  Lemma~\ref{lem:quotient-embedding} gives
  \[
    \Phi_1(f)
    =\pi_{I'^1}(F^*f)
    =j'\bigl(\pi_{I'}(F^*f)\bigr)
    =\bigl(j'\circ\mathcal{E}_{\mathrm{ae}}[F]\bigr)(f) .
  \]
  That is, $\Phi_1|_L=j'\circ\mathcal{E}_{\mathrm{ae}}[F]=\Phi_2|_L$.

  The system $(X',L'^1,I'^1)$ is a Daniell closed integration system
  (Lemma~\ref{lem:daniell-ext-is-complete}), and
  Lemma~\ref{lem:quotient-is-banach} shows that $\widetilde{L'^1}$ is a Banach
  space, so the \textbf{uniqueness} in
  Theorem~\ref{thm:operator-extension} applies and we obtain
  $\Phi_1=\Phi_2$.
\end{proof}

\begin{remark}
  Note that this proof relies on the uniqueness of
  Section~\ref{chap:operator}.
  Within the framework of Section~\ref{chap:functoriality} alone there is no
  means of saying that the two extensions agree, because they are morphisms of
  different categories.
\end{remark}

\subsubsection{Relation between the two main theorems}\label{sec:two-chapters}

\begin{corollary}\label{cor:ch3-uniqueness-from-ch4}
  Let $A$ be an object of $\mathbf{Int}_{\mathrm{ae}}$, let $B$ be an object of
  $\mathbf{DInt}_{\mathrm{ae}}$, and let
  $[H],[H']:\mathcal{D}_{\mathrm{ae}}A\to B$ be morphisms.
  Then
  \[
    \iota_{D,\mathrm{ae}}([H])\circ[i_A]
    =\iota_{D,\mathrm{ae}}([H'])\circ[i_A]
    \quad\Longrightarrow\quad
    [H]=[H']
  \]
  holds. That is, the uniqueness part of
  Theorem~\ref{thm:main-ae} follows from
  Theorem~\ref{thm:operator-extension} of
  Section~\ref{chap:operator}.
\end{corollary}
\begin{proof}
  Apply $\mathcal{E}_{\mathrm{ae}}$ to both sides.
  By Lemma~\ref{lem:E-descends-to-ae} the assignment
  $\mathcal{E}_{\mathrm{ae}}$ is a functor, and
  $\mathcal{E}_{\mathrm{ae}}[i_A]$ is the morphism determined by the inclusion
  $L\subset L^1$, so $\mathcal{E}_{\mathrm{ae}}[H]$ and
  $\mathcal{E}_{\mathrm{ae}}[H']$ are operator morphisms into the quotient
  $L^1\to\widetilde{M}$ agreeing on $L$.
  Since $B$ is Daniell closed, $\widetilde{M}$ is a Banach space by
  Lemma~\ref{lem:quotient-is-banach}, and the uniqueness in
  Theorem~\ref{thm:operator-extension} gives
  $\mathcal{E}_{\mathrm{ae}}[H]=\mathcal{E}_{\mathrm{ae}}[H']$.
  As $\mathcal{E}_{\mathrm{ae}}$ is faithful
  (Lemma~\ref{lem:E-descends-to-ae}), we obtain $[H]=[H']$.
\end{proof}

\begin{remark}[relation between the two sections]\label{rem:two-chapters-relation}
  By contrast, the \textbf{existence} part of
  Theorem~\ref{thm:main-ae} does not follow from the theorem of
  Section~\ref{chap:operator}.
  Given $[F]:A\to\iota_{D,\mathrm{ae}}B$,
  Theorem~\ref{thm:operator-extension} supplies an extension of
  $\mathcal{E}_{\mathrm{ae}}[F]$ as a morphism of
  $\mathbf{Int}_{\mathrm{op},\mathrm{ae}}$, but it does not guarantee that this
  lies in the image of $\mathcal{E}_{\mathrm{ae}}$, that is, that it comes from
  a map of base spaces, because $\mathcal{E}_{\mathrm{ae}}$ is not full
  (Remark~\ref{rem:E-ae-faithful-not-full}).
  To obtain existence one needs the construction specific to
  Section~\ref{chap:functoriality}
  (Lemma~\ref{lem:daniell-ext-is-functor-1}), by which the extension is
  obtained without changing the map of base spaces.

  The contributions of the two theorems divide as follows.

  \begin{center}
  \begin{tabular}{@{}lll@{}}
    \toprule
     & Existence & Uniqueness \\
    \midrule
    Section~\ref{chap:functoriality}, non-ae & needs its own construction & follows automatically (no content) \\
    Section~\ref{chap:functoriality}, ae & as above & by density \\
    Section~\ref{chap:operator} (operators) & not obtainable from the operator version & by contractivity and density \\
    \bottomrule
  \end{tabular}
  \end{center}

  Care is needed in how this is phrased.
  Since $\mathcal{E}_{\mathrm{ae}}$ is faithful,
  $\mathbf{Int}_{\mathrm{ae}}$ is embedded (without changing objects) in
  $\mathbf{Int}_{\mathrm{op},\mathrm{ae}}$.
  In that sense \textbf{the category of
  Section~\ref{chap:functoriality} is a wide subcategory of the category of
  Section~\ref{chap:operator}}, and it is not accurate to say that ``neither is
  a special case of the other''.
  Nor is Section~\ref{chap:functoriality} in the relation of completing a single
  theorem together with Section~\ref{chap:operator}: it proves both existence
  and uniqueness on its own.
  The relations that can be stated precisely are the following.

  \begin{itemize}
    \item $\mathcal{E}_{\mathrm{ae}}$ is faithful but not full
      (Remark~\ref{rem:E-ae-faithful-not-full}).
    \item The two reflections are compatible under this functor
      (Proposition~\ref{prop:E-commutes-with-D}).
    \item The \emph{uniqueness} of the base-space-map version can be derived
      from the operator version
      (Corollary~\ref{cor:ch3-uniqueness-from-ch4}).
    \item But \textbf{that the extension stays in the image of
      $\mathcal{E}_{\mathrm{ae}}$, that is, that it comes from a map of base
      spaces, cannot be derived from the operator version alone.}
    \item Furthermore, \textbf{$\mathbf{Int}_{\mathrm{op},\mathrm{ae}}$ sees
      neither the base space nor the realization as a function space}
      (Proposition~\ref{prop:Q-full-faithful},
      Corollary~\ref{cor:op-ae-iso-criterion}).
      Objects are distinguished only up to the pre-AL space
      $(\widetilde{L},\widetilde{I})$, and what the Daniell extension gives is
      its Banach completion
      (Proposition~\ref{prop:D-op-is-banach-completion}).
  \end{itemize}

  The last two points are the reasons for keeping
  Section~\ref{chap:functoriality} as an independent section.
  The first is a reason at the level of theorems and the second at the level of
  categories, and the latter is the stronger of the two.
  Only the side of Section~\ref{chap:functoriality} retains the distinction
  between an abstract completion and a concrete realization
  (Subsection~\ref{sec:prior-work}) as categorical structure.
  For details see Subsection~\ref{sec:what-op-ae-sees}.
\end{remark}

